% Ч1.3 — Проценты: процент от числа, число по проценту, процентное отношение
% Старая методичка: стр. 91–92 (старое задание 10)
% Источники: методичка по заданию 13 (часть 1), методичка по заданию 11 (проценты и доли).
% Стиль — стиль-изложения.md

\begin{tema}{Ч1.3}
\opr{Процент} — это сотая часть числа: $1\%=\frac{1}{100}$.
Чтобы найти $p\%$ от числа, нужно это число умножить на $\frac{p}{100}$.
Перед решением важно понять, \emph{что принято за 100\,\%}, — от этого зависит всё остальное.

\begin{pfTablica}{lc}
Что найти & Как\\
$p\%$ от числа $x$ & $x\cdot\dfrac{p}{100}$\\
Число, если $p\%$ от него равны $a$ & $a\cdot\dfrac{100}{p}$\\
Сколько процентов $a$ составляет от $b$ & $\dfrac{a}{b}\cdot100\%$\\
На сколько процентов $A$ больше $B$ & $\dfrac{A-B}{B}\cdot100\%$\\
На сколько процентов $B$ меньше $A$ & $\dfrac{A-B}{A}\cdot100\%$\\
\end{pfTablica}

\opr{Деление в отношении.} Чтобы разделить число $S$ в отношении $a:b$,
разделите его на $a+b$ равных частей: первая доля равна $S\cdot\frac{a}{a+b}$,
вторая — $S\cdot\frac{b}{a+b}$. Например, 120 в отношении $2:3$ — это 48 и 72.

\begin{primer}
У Веры было 1200 рублей. 18\,\% этой суммы она потратила на тетрадь.
Сколько стоила тетрадь?
Решение: $\frac{18}{100}\cdot1200=216$. Ответ: 216 рублей.
\end{primer}

\begin{oshibka}
«$A$ больше $B$ на $p\%$» и «$B$ меньше $A$ на $p\%$» — разные утверждения:
50 больше 40 на 25\,\%, а 40 меньше 50 на 20\,\%. За 100\,\% берётся то,
\emph{с чем} сравнивают.
\end{oshibka}
\end{tema}
